\documentclass[10pt,a4paper]{amsart}
\usepackage[margin=19mm]{geometry}
\usepackage{amsmath,amssymb,mathtools}
\usepackage[scr=rsfs,cal=euler]{mathalfa}
\usepackage{xfrac}
\usepackage[unicode,hidelinks,bookmarks=false]{hyperref}
\newcommand{\ie}{i.e.\ }
\newcommand{\cf}{cf.\ }
\newcommand{\resp}{respectively\ }
\newcommand{\Subsection}{\subsection}
\providecommand{\nobreakdash}{\nobreak\hbox{-}}
\setlength{\emergencystretch}{2em}
\multlinegap0pt
\usepackage{braket}
\usepackage{amssymb}
\usepackage{multirow}
\usepackage[shortlabels]{enumitem}
\usepackage{tikz}
\usepackage{mathtools}
\newtheorem{lemma}{Lemma}
\newcommand{\aditya}[1]{\marginpar{+}{\bf Aditya's remark}: {\em #1}}
\newcommand{\bbC}{\mathbb{C}}
\newcommand{\bfX}{\mathbf{X}}
\newcommand{\bfY}{\mathbf{Y}}
\newcommand{\bfx}{\mathbf{x}}
\newcommand{\bfy}{\mathbf{y}}
\newcommand{\bfz}{\mathbf{z}}
\newcommand{\calX}{\mathcal{X}}
\newcommand{\calY}{\mathcal{Y}}
\newcommand{\calZ}{\mathcal{Z}}
\newcommand{\meng}[1]{\marginpar{+}{\bf Ruoyu's remark}: {\em #1}}
\title{Quantum advantage in zero-error function computation with side information}
\author{Ruoyu Meng, Aditya Ramamoorthy}
\date{Source: arXiv:2402.01549v4 (CC BY 4.0)}
\begin{document}
\maketitle

\noindent\emph{Source and licence.} Quantum advantage in zero-error function computation with side information. Version arXiv:2402.01549v4 (CC BY 4.0), distributed by arXiv under the Creative Commons Attribution 4.0 International licence, http://creativecommons.org/licenses/by/4.0/, as recorded in the arXiv licence metadata for this exact version and shown on the abstract page https://arxiv.org/abs/2402.01549v4. Full licence notice: \texttt{licenses/arxiv-2402.01549-CC-BY-4.0.txt}.\par\smallskip

\noindent\textit{Editorial scope.} Counted unit: the article's lemma lemma:lemma4 (source0.tex lines 1346--1348). The article's complete original proof of it is delivered with it (source0.tex lines 1349--1374, one \texttt{proof} environment of 26 lines). No mathematics is written, repaired or re-derived: the statement and the proof are the article's own lines, in the article's own notation, and no retained line is substituted or re-flowed.  No further source-local context is needed: every cross-reference of every retained line resolves inside the two delivered spans.  Nothing was omitted from any retained span, so the package carries no omission record.  No retained line contains a citation, so the wrapper carries no bibliography and no bibliography entry.  No retained line contains a figure environment, an image-inclusion command or drawing code, so no image is delivered and the package declares no asset dependency.  Editorial formatting only, in addition: an AMS \texttt{amsart} preamble that restates the article's own declarations and macros used by the retained lines, namely the environments \texttt{lemma} on separate counters, together with the standard packages those lines need.  The retained environments print in this document as Lemma 1 in order of appearance, whereas the article prints the same items with the numbering of its own full document.  The complete native source of the article as distributed in the arXiv e-print for this exact version is bundled unchanged as \texttt{sources/153945/source0.tex}.
\begin{lemma}\label{lemma:lemma4}
For a function $f:\calX\times\calY \mapsto\calZ$ with p.m.f. $p_{XY}(\cdot,\cdot)$, let $G$ denote the $f$-confusion graph. The optimal quantum rate for computation over $m$ instances is $\frac{\log_2\xi(\overline{G^{(m)}})}{m}$ ($\overline{G^{(m)}}$ denotes the complement of the $f$-confusion graph over $m$-instances $G^{(m)}$).
\end{lemma}

\begin{proof}
An orthogonal representation $\phi:\overline{G^{(m)}}\mapsto \bbC^{\xi(\overline{G^{(m)}})}$  induces a quantum code as follows. Suppose Alice and Bob get  $\bfx $ and $\bfy$ respectively.  Alice sends $\ket{\phi(\bfx)}\bra{\phi(\bfx)}$ to Bob.  Bob chooses his POVM to be $\{\Pi^{\bfz,\bfy}  :\bfz\in \calZ^m\}$ where $\Pi^{\bfz,\bfy}$ is the projector onto $S_{\bfz,\bfy}:=\text{span}\{\ket{\phi(\bfx)}:p_{\bfX\bfY}(\bfx,\bfy)>0, f^{(m)}(\bfx,\bfy) 
 = \bfz \}$. 
 If $\sum_{\bfz}\Pi^{\bfz,\bfy}\ne I$, Bob adds $I-\sum_{\bfz}\Pi^{\bfz,\bfy}$ to complete a POVM. If $\bfx,\bfx'$ is such that $p_{\bfX\bfY}(\bfx,\bfy)p_{\bfX\bfY}(\bfx',\bfy)>0$, and $f^{(m)}(\bfx,\bfy)=\bfz,f^{(m)}(\bfx',\bfy)=\bfz'$ for some $\bfz\ne \bfz'$, then $\ket{\phi(\bfx)} \in S_{\bfz,\bfy}$ and $\ket{\phi(\bfx')}\in S_{\bfz',\bfy}$. Since $\phi$ is an orthogonal representation of $\overline{G}$, then  $\Pi^{\bfz,\bfy}\perp\Pi^{\bfz',\bfy}$, which implies that  the code is zero-error.   The code has rate $\frac{1}{m}\log_2\xi(\overline{G^{(m)}})$ as the orthogonal representation $\phi$ is $\xi(\overline{G^{(m)}})$-dimensional.

%Now, we show the protocol is zero-error. Let $\bfx$ and $\bfy$ be the same as the  paragraph above.  Denote $\bfz  = f^{(m)}(\bfx,\bfy)$. Since $\ket{\phi(\bfx)}\in S_{\bfz,\bfy}$ and $\ket{\phi(\bfx)}\notin S_{\bfz',\bfy}$ for all $\bfz\ne \bfz'$, Bob's measurement result will be $\bfz$ with probability $1$.   
%\aditya{Is this para needed?} \meng{The paragraph is obsorbed into the previous paragraph}

%\meng{One reviewer mentioned that the braket notation is not defined, but I feel it is well-known for people in quantum information.} 
Conversely, a code   of rate $\frac{1}{m}\log_2d$ induces an orthogonal representation $\psi:V(G^{(m)})\mapsto \bbC^{d}$. Let $\rho_{\bfx}$ be a $d\times d$  density operator associated with input $\bfx$. If $\rho_{\bfx}$ is a pure state $\ket{v}\bra{v}$ for some $\ket{v}\in \bbC^d$, then we set  
 $\psi(\bfx) = \ket{v}$. If $\rho_{\bfx}$ is a mixed state with its density operator written as  
 $\rho_{\bfx} =\sum_{i=1}^k \lambda_i\ket{v_i}\bra{v_i}\text{ with }\lambda_i\ne 0\text{ for }i=1,\dots,k.$
Then let $l\in[k]$ be arbitrary and $\psi(\bfx ) = \ket{v_l}$. 

%\meng{I modified the next paragraph such that a reviewer not knowing QIT can understand from the perfectly distinguishable states that I just added in section I-A quantum setting.}


Now we check it is indeed an orthogonal representation. Let    $(\bfx,\bfx')\in E(G^{(m)})$ be such that $p_{\bfX\bfY}(\bfx,\bfy)p_{\bfX\bfY}(\bfx',\bfy)>0$ and $f^{(m)}(\bfx,\bfy)\ne f^{(m)}(\bfx',\bfy)$.  This code  is zero-error, so the POVM $\{\Lambda^\bfz_{f(\cdot,\cdot), \bfy}\}_{\bfz \in \calZ^m}$ must distinguish $\rho_{\bfx},\rho_{\bfx'}$ perfectly. This implies that $\rho_{\bfx}\perp\rho_{\bfx'}$. We have 
%Now we check it is indeed an orthogonal representation. Let    $(\bfx,\bfx')\in E(G^{(m)})$.  This protocol  is zero-error, so $\rho_{\bfx} =\sum_{i=1}^k \lambda_i\ket{v_i}\bra{v_i},\rho_{\bfx'} =\sum_{i=1}^k \lambda_i'\ket{v_i'}\bra{v_i'}$ must be perfectly distinguishable, i.e. $\rho_{\bfx}\perp\rho_{\bfx'}$. We have 
\begin{align*}
   &  \sum_{i=1}^k \lambda_i\ket{v_i}\bra{v_i}  \perp  \sum_{i=1}^k \lambda_i'\ket{v_i'}\bra{v_i'}  \Rightarrow \ket{v_i}\perp\ket{v_j'}\text{ for all }i,j
%   \\& 
   \Rightarrow \psi(\bfx) \perp \psi(\bfx').
\end{align*}
Since $(\bfx,\bfx')\in E(G)$ is arbitrary, it implies that $\psi$ is an orthogonal representation.  
\end{proof}

\end{document}
